\documentclass{situationist-preprint}

\title{Obligations to No One in Particular}
\msid{ETH-08-a}
\msfamily{ETH-08}
\msdiscipline{Normative and Applied Ethics}
\mstype{Principal manuscript}
\mskeywords{future generations; non-identity problem; person-affecting principles; impersonal obligation; directed duties; Parfit; intergenerational justice; climate}

\begin{document}
\maketitle

\begin{abstract}
Do we owe anything to future generations? We argue that we do, but that the obligations are not owed to anyone in particular: they are impersonal. The argument turns on the non-identity problem. Policies of the scale that concern future generations, such as the depletion of resources or the emission of greenhouse gases, affect which people will exist, and those who exist as a result of a harmful policy are not made worse off by it, since they would not otherwise have existed. We show that they are not wronged either: a directed claim can ground a wrong only if there is an alternative open to the agent in which its holder exists and the claim is honored, and for the people born under such a policy there is none. Yet the policies are wrong. The obligation they violate therefore has no particular claimant. We review the main attempts to find one, through rights, non-comparative harm, contractualist types and de dicto interests, and show that each either fails in a central case or concedes that the obligation is owed to no one in particular. The obligation is grounded in the impersonal value of outcomes.
\end{abstract}

\section{The problem}

Many of our choices will affect people who do not yet exist. A society that consumes its forests and fisheries, that stores nuclear waste where it will leak in five centuries, or that warms the climate over the next hundred years, will leave its successors worse conditions than it might have. It is natural to think that such a society wrongs its successors, that it fails in what it owes them, and that they will have a complaint against it. The view is shared by many who disagree about almost everything else in ethics, and it informs both the ordinary language of intergenerational justice and the reasoning of international law.

\citet[ch.~16]{Parfit1984} showed that the natural thought faces a difficulty, which he called the \emph{non-identity problem}. Who exists in the future depends on which people meet and when they conceive their children. A policy that affects the conditions of life of a whole society over decades changes these facts for almost everyone, and after a generation or two almost everyone who exists under the policy is someone who would not have existed under its alternative. Consider Parfit's case of \emph{Depletion}. A community must choose between conserving a resource and depleting it. If it depletes, the quality of life over the next two centuries will be slightly higher than it would have been under conservation, but thereafter it will be much lower, though still worth living. The people who live in three centuries under depletion would not have existed under conservation. Depletion is not worse for them, since the alternative is not that they would have had better lives but that they would not have existed. It seems that depletion is worse for no one. And yet it seems clearly wrong.

The problem had been noticed before Parfit, in different forms, by \citet{Schwartz1978}, \citet{Adams1979} and \citet{Kavka1982}. Schwartz concluded that we have no obligation to provide for the distant future; Adams used the point in theodicy, arguing that God does not wrong us by creating a world less good than another in which we would not exist. Parfit's contribution was to show that the problem cannot be dismissed by accepting Schwartz's conclusion, since that conclusion is incredible, and cannot be solved by any principle that appeals to what is worse for particular people. He concluded that the solution must lie in an impersonal principle, which he hoped would also solve the problems of different-number choices, and which he did not find.

We defend Parfit's diagnosis in a stronger form. We argue not only that depletion is worse for no one, but that it wrongs no one: there is no particular person to whom the obligation to conserve is owed. The obligation exists; it is impersonal. Section~\ref{s:framework} sets out the framework. Section~\ref{s:argument} gives the argument. Section~\ref{s:ground} identifies the ground of the obligation and connects it with the population axiology of ETH-09. Section~\ref{s:history} reviews the attempts to find a particular claimant and the cases that defeat each. Section~\ref{s:cases} applies the view to further cases. Section~\ref{s:objections} answers objections, and Section~\ref{s:limits} records what the argument does not establish.

\section{Framework}\label{s:framework}

\begin{definition}[Directed and impersonal obligation]\label{d:directed}
An obligation is \emph{directed}, or \emph{owed to} a person, if there is a particular person such that failing to fulfil it wrongs that person. An obligation is \emph{impersonal} if failing to fulfil it is wrong but wrongs no particular person.
\end{definition}

The phrase ``particular person'' is to be read \emph{de re}. An obligation that is owed to ``whoever will live in the twenty-third century'' is owed to a particular person only if there is someone of whom it is true that the obligation is owed to her. The description ``whoever will live then'' does not, by itself, pick out such a person, since which people will satisfy it depends on what we do. Whether an obligation is owed to someone de re is the question at issue.

\begin{definition}[Worse for]\label{d:worse}
An outcome $O$ is \emph{worse for} a person $p$ than an outcome $O'$ if $p$ exists in both and her life goes less well in $O$ than in $O'$.
\end{definition}

Definition~\ref{d:worse} is the comparative notion of being made worse off. It requires that $p$ exist in both outcomes. We do not assume that coming into existence can be better or worse for someone than not existing; we assume only that it is not \emph{worse} for someone whose life is worth living to exist than not to exist, which is common ground.

\begin{definition}[Identity-affecting choice]\label{d:ia}
A choice between options is \emph{identity-affecting} with respect to a person $p$ who exists if one option is chosen, if $p$ would not have existed had another option been chosen.
\end{definition}

\begin{assumption}[Origin dependence]\label{a:origin}
A person who in fact develops from a particular pair of gametes could not have existed without developing from that pair.
\end{assumption}

Assumption~\ref{a:origin} is a version of the necessity of origin \citep[lecture III]{Kripke1980}. The argument needs only a weak form of it: that a person conceived months later, from different gametes, by different parents, would be a different person. It does not need the strong form, which forbids any variation in origin.

\begin{principle}[Comparative condition on harming]\label{p:harm}
A choice harms a person only if it makes her worse off, in the sense of Definition~\ref{d:worse}, than she would otherwise have been.
\end{principle}

\begin{principle}[Claim condition on wronging]\label{p:claim}
A choice wrongs a person only if it harms her or infringes a claim she holds against the agent; and it infringes a claim held by $p$ only if some alternative open to the agent is one in which $p$ exists and the claim is honored.
\end{principle}

Principle~\ref{p:claim} has two parts. The first allows that one may wrong without harming: Woodward's passenger, refused a seat on an aircraft because of his race, is wronged even if the aircraft crashes and the refusal saves his life \citep{Woodward1986}. The second part says what a claim is. A claim is a claim to be treated in a certain way, and its infringement consists in not treating its holder that way when one could have. In Woodward's case, the airline could have seated the passenger, and its refusal infringed his claim. A claim of $p$ that no option open to the agent could honor, because $p$ exists only in options in which it is not honored, is not a claim against the agent: there is no way the agent could have treated $p$ that the claim calls for. We defend the second part further in Section~\ref{s:objections}.

\begin{principle}[Depletion]\label{p:depletion}
In Depletion, it is wrong for the community to choose depletion.
\end{principle}

Principle~\ref{p:depletion} is the datum. Few deny it, and those who do, as we will see, accept a cost that most regard as disqualifying.

\section{The argument}\label{s:argument}

\begin{lemma}\label{l:ia}
The choice between depletion and conservation is identity-affecting with respect to almost everyone who exists, under either option, more than a generation after the choice.
\end{lemma}

\begin{proof}
A policy affecting the economy and the conditions of life of a whole society changes, for nearly everyone, the details of daily life: where people work and live, whom they meet, and when they conceive. A difference of a few minutes in the timing of conception results in a different pair of gametes, and by Assumption~\ref{a:origin}, in a different person. Within a generation, almost everyone born differs from those who would have been born under the alternative; and since the people of each later generation are the children of the people of earlier ones, the difference propagates. After a few generations no one would exist under both options.
\end{proof}

\begin{lemma}\label{l:noharm}
No person who exists under depletion more than a generation after the choice is harmed by it.
\end{lemma}

\begin{proof}
By Lemma~\ref{l:ia}, such a person would not have existed under conservation. By Definition~\ref{d:worse}, depletion is not worse for her than conservation, since she does not exist in the latter. By hypothesis her life is worth living, so depletion is not worse for her than her non-existence would have been. By Principle~\ref{p:harm}, she is not harmed.
\end{proof}

\begin{lemma}\label{l:noclaim}
No person who exists under depletion more than a generation after the choice has a claim against the community that depletion infringes.
\end{lemma}

\begin{proof}
Suppose $p$ exists under depletion and has a claim against the community, for instance to a certain standard of resources. By Lemma~\ref{l:ia}, the only option open to the community in which $p$ exists is depletion, and under depletion the claim is not honored. By the second part of Principle~\ref{p:claim}, depletion does not infringe the claim.
\end{proof}

\begin{theorem}\label{t:main}
Depletion wrongs no person who exists under it more than a generation after the choice. It also wrongs no person who exists under conservation, and no one who exists under neither.
\end{theorem}

\begin{proof}
By Principle~\ref{p:claim}, a choice wrongs someone only if it harms her or infringes a claim of hers. For people who exist under depletion more than a generation later, Lemmas~\ref{l:noharm} and \ref{l:noclaim} exclude both. A person who exists only under conservation does not exist if depletion is chosen; she is not made worse off (Definition~\ref{d:worse} requires existence in both outcomes), and since there is no one who has the claim, no claim is infringed. A person who exists under neither option is not affected by the choice.
\end{proof}

The theorem concerns the distant future. People who exist under both options, those already alive and those born in the first years after the choice, are affected in the ordinary way, and in Parfit's version of the case depletion benefits them.

\begin{corollary}\label{c:impersonal}
The obligation not to choose depletion is impersonal.
\end{corollary}

\begin{proof}
By Principle~\ref{p:depletion}, choosing depletion is wrong. By Theorem~\ref{t:main}, the wrongness does not consist in wronging any of the people who will exist more than a generation after the choice, and since depletion benefits those who exist under both options, it does not consist in wronging them. So it wrongs no particular person, and by Definition~\ref{d:directed} the obligation is impersonal.
\end{proof}

The corollary generalizes. The policies that concern future generations as such, policies about resources, the environment, the stock of knowledge and institutions, and population, are all identity-affecting at the scale of generations. Where they would leave future people worse conditions than an alternative, and those people have lives worth living, the argument applies. We conclude that our obligations to future generations are impersonal: we owe it to no one in particular to leave a good world behind.

\section{The ground of the obligation}\label{s:ground}

If the obligation is not owed to anyone, what grounds it? Parfit's answer, in same-number cases, was his principle Q: if in either of two outcomes the same number of people would ever live, it would be worse if those who live are worse off, or have a lower quality of life, than those who would have lived \citep[Part Four]{Parfit1984}. The principle is impersonal: it compares the quality of life in two outcomes without asking whether the people are the same. Parfit held, and we agree, that the non-identity problem makes no moral difference in same-number cases. The choice of depletion is wrong for the same reason, and to the same degree, as it would be if the same people would exist under both options. This is Parfit's \emph{No-Difference View}.

Principle Q does not extend to different-number cases, and the theory that Parfit hoped would extend it, Theory X, he never found. ETH-09 supplies it. On the population axiology defended there, the value of an outcome is impersonal and additive across lives, so that one outcome is better than another if it contains more value summed over all the lives in it. In same-number cases, this coincides with principle Q. In different-number cases, it implies the repugnant conclusion, which ETH-09 argues is true and, correctly understood, not repugnant. The impersonal obligations of the present manuscript are obligations to bring about, or not to forgo, outcomes that are better in this sense, within limits set by the deontic principles discussed below. ETH-20 relies on the same axiology when it argues that the absence of the goods a life would contain, though bad for no one, is not neutral.

\begin{proposition}\label{p:ground}
If the value of outcomes is impersonal in the sense of ETH-09, then the wrongness of depletion is explained without appeal to any wronged person, and the No-Difference View holds.
\end{proposition}

\begin{proof}
On the impersonal axiology, the value of an outcome depends on the values of the lives in it and not on who lives them. Depletion produces an outcome whose lives, after two centuries, are much less good than those of conservation, with the same number of people, and so produces a worse outcome. The obligation not to choose a worse outcome when a better one is available at moderate cost explains why depletion is wrong. The explanation does not mention the identities of the people affected, so it gives the same verdict whether or not the choice is identity-affecting.
\end{proof}

Two features of the obligation follow. First, impersonal obligations are not obligations to maximize. The step from the value of outcomes to what agents ought to do is governed by deontic principles: the duty of beneficence is moderate (ETH-06), and the combined duties of beneficence and rescue leave room for a personal prerogative (ETH-28). The obligation to future generations is an instance of the duty not to make outcomes worse when the cost of avoiding it is moderate. Collective agents such as states are bound by it as individuals are, and the scale at which they act makes the cost of compliance, relative to the stakes, small.

Second, the obligation has a definite content even though it has no particular claimant. It requires that we leave the future no worse, in quality of life, than we reasonably can; and since the population axiology of ETH-09 counts the number of lives as well as their quality, it also counts against outcomes in which future populations are much smaller than they could have been with no loss in quality. It does not require that we bring about the largest possible population, since the deontic principles and the autonomy of procreative decisions intervene.

\section{The attempts to find a claimant}\label{s:history}

The literature since Parfit has largely consisted of attempts to show that the future people are, after all, wronged, so that the obligation is directed. We review the principal attempts and the case that defeats each.

\paragraph{Narveson: no obligation.} Before Parfit, \citet{Narveson1967} argued that utilitarian reasons are reasons to make people happy, not to make happy people: there is no reason to bring someone into existence for her own sake. Narveson's view is person-affecting, and applied to Depletion it implies, with Schwartz, that there is no obligation to conserve, since no one is made worse off. Depletion is the case that defeats it.

\paragraph{Bullet-biting.} \citet{Boonin2014} argues at length that the non-identity problem has no solution and that we should accept its conclusion: in the cases in which a choice makes no one worse off, it is not wrong on the grounds usually given. Boonin's argument is that every attempt to resist the conclusion relies on a premise more doubtful than the denial of the conclusion. We agree with Boonin about the attempts to find a wronged person, which are discussed below. We disagree with his conclusion because he treats the five premises of the argument as the only premises available, and so overlooks the possibility that the wrongness is impersonal. Theorem~\ref{t:main} is a version of his conclusion about wronging; Corollary~\ref{c:impersonal} draws a different conclusion about wrongness.

\paragraph{Rights.} \citet{Woodward1986} argued that one may wrong a person by violating her rights even if one does not make her worse off, and that future people have rights, for example to an environment of a certain quality, which depletion violates. His airline case shows that wronging does not require harming, and we accept this in the first part of Principle~\ref{p:claim}. But the right appealed to in Depletion has a peculiar structure. Unlike the passenger's right, it is a right that cannot be honored without its holder's not existing. If the community conserves, the rights-holder does not exist; if it depletes, the right is violated. The right therefore cannot be a reason for the community to conserve \emph{for the sake of} its holder. The people in Depletion, if asked, would rationally waive any such right if waiving it were a condition of their existence. The second part of Principle~\ref{p:claim} states the point generally.

\paragraph{Non-comparative harm.} \citet{Shiffrin1999} and \citet{Harman2004} argued that harm need not be comparative: to cause someone pain, early death, bodily damage or deformity is to harm her, whether or not she would otherwise have been better off. On this view, a choice can harm people whom it brings into existence. Harman uses the view to explain cases such as Parfit's fourteen-year-old girl, who chooses to have a child now rather than later, when she could give it a better start; the child she has now suffers harms of the non-comparative kind, and her reasons against causing them are not outweighed by any benefit to the child, since the benefit of existence does not count. The view handles cases in which future people suffer harms of a recognizable kind. It does not handle Depletion, in which people under depletion have lives that are less good than lives under conservation but contain no non-comparative harm beyond what lives under conservation contain. They are not injured, diseased or deprived; there is simply less of the good in their lives. If non-comparative harm extended to having less of the good, it would include almost every life, since almost every life contains less of the good than some possible life, and Harman's view would imply that almost every act of procreation harms the child.

\paragraph{Hanser.} \citet{Hanser1990} argued that a person can be harmed by an act if it causes her to suffer a harmful event, such as an injury, even if she would not have existed without the act. The account is a non-comparative account of harming events rather than states, and it shares the limitation of the account of Shiffrin and Harman: it applies to cases with harmful events, and Depletion need contain none.

\paragraph{Contractualist types.} \citet{Kumar2003} argued, within contractualism, that a person can be wronged in virtue of the \emph{type} of person she is. The child of the fourteen-year-old girl is wronged because people of the type ``child of a mother who could reasonably have given it a better start'' have a claim against being treated as the girl treated hers. The proposal moves the claim from the particular person to a type that the person occupies, and it is the type, not the person, that the obligation tracks. We regard this as a concession. If the obligation is owed to occupants of a type, and the identity of the occupant makes no difference to the obligation, then the obligation is not owed to anyone de re. Kumar's view correctly identifies the content of the obligation, but in our terms the obligation it identifies is impersonal.

\paragraph{De dicto interests.} \citet{Hare2007} argued that we should care about the interests of future people \emph{de dicto}: a parent should want ``my child'', whoever that is, to be healthy, and a society should want ``the people of the twenty-third century'' to flourish. On this view, depletion is bad for the people of the twenty-third century de dicto, though for none of them de re. We accept the description and draw the same conclusion as with Kumar. An obligation grounded in de dicto interests is one that holds whoever the people are; there is no particular person to whom it is owed. Hare's view is a way of expressing the impersonal obligation in the vocabulary of interests.

\paragraph{Variabilism.} \citet{Roberts2010} defends a person-affecting view on which the moral significance of a loss depends on whether the person who incurs it exists in the outcome. The view is designed to explain why we need not create happy people while holding that we must not create people with worse lives than we could. It handles the non-identity cases by treating the loss in the depleted outcome as a loss incurred by existing people relative to what ``the people in that position'' would have had. As with Kumar and Hare, the comparison is between positions, not persons, and the view concedes that no particular person is made worse off.

\paragraph{Non-moral creation.} \citet{Heyd1992} argued that decisions about whom to create lie outside the scope of morality, since morality concerns the interests of existing people. The view avoids the non-identity problem by denying that Depletion's effects on the distant future are a moral matter, and it is defeated by Depletion for the same reason as Narveson's.

The attempts divide into two groups. Those that look for a particular wronged person fail in Depletion, because the people in Depletion are neither harmed in any sense that applies to the case nor holders of claims that the community could honor. Those that succeed do so by relocating the claim in a type, a position or a de dicto description. They agree with us about the content of the obligation and disagree only in describing it as directed. Since the claimant they identify is no one in particular, the disagreement is verbal.

\section{Further cases}\label{s:cases}

\paragraph{Risky Policy.} \citet[Part Four]{Parfit1984} describes a community choosing between two energy policies, one of which carries a risk that, centuries later, an earthquake will release radiation and kill thousands. The people killed would not have existed under the safer policy. Unlike Depletion, Risky Policy involves a harmful event, and Hanser's account applies. On our view, the deaths are a loss that the impersonal axiology counts, and the choice of the risky policy is wrong if the expected loss is not outweighed by the gains. Whether the people killed are also wronged, in addition, by a harmful event, we need not decide. If they are, it is because the event of being killed by radiation is a harm in Hanser's sense, and that does not depend on the identity-affecting character of the choice. The intergenerational obligation not to adopt the risky policy is impersonal, and it holds whether or not the policy also causes harmful events to particular future persons.

\paragraph{The fourteen-year-old girl.} The girl's child has a life worth living. Had she waited, she would have had a different child with a better start. The choice is identity-affecting, and the child is not worse off than he would otherwise have been. On our view, the girl acts wrongly, if she does, because she brings about a worse outcome when a better was available at moderate cost: the cost to her of waiting is the cost against which the impersonal reason is weighed. The verdict matches common opinion in being mild. Since the cost to her may be considerable, and procreative autonomy is a component of the good, the impersonal reason may not yield an obligation. ETH-08-b discusses this and related procreative cases.

\paragraph{Two medical programmes.} Parfit compares two programmes, each of which would prevent a condition that impairs the lives of a thousand children a year. The first tests pregnant women and treats the condition in the womb; the second tests women before conception and advises those who are affected to delay conceiving. The first programme prevents the impairment in children who will exist anyway; the second prevents the existence of the impaired children by causing different children to exist. If only one programme can be funded, and they would prevent the same number of impaired lives, it seems to make no moral difference which. A person-affecting view favors the first, since only the first benefits particular children. The impersonal view gives the two equal weight, in accordance with the No-Difference View, and we take this to be correct.

\paragraph{Climate.} Climate change is the case that gives the problem its current urgency \citep{Gardiner2011,Broome2012}. Emissions over the next century will affect the conditions of life of billions of people born after it, all of whom, by Lemma~\ref{l:ia}, would not have existed under a different emissions policy. The obligation to reduce emissions, insofar as it concerns those people, is impersonal. This has two consequences for climate ethics. First, arguments that rest on the rights or claims of particular future people, such as a right to a stable climate, must be restated in impersonal terms; restated, they lose nothing of their force, since the impersonal obligation has the same content. Second, the discounting of future well-being cannot be justified by the claim that future people, being unidentified, have weaker claims on us: on the impersonal view, the value of a life does not depend on when it is lived or on whether its subject can be identified now. \citet{Broome2012} reaches the same conclusion about pure time discounting by a different route.

\paragraph{Buried glass.} Someone leaves broken glass in undergrowth, and a century later a child is cut by it. The choice to leave the glass is not identity-affecting, since leaving it does not change who exists. The child is harmed in the ordinary comparative sense and is wronged. This case shows the limits of our thesis. Not every obligation concerning people who will live in the future is impersonal; obligations not to harm people who will exist whatever we do are directed, even if we cannot now identify their holders. What is impersonal is the class of obligations to future generations as such: those concerning the conditions of life at the scale of generations, which are discharged by policies that affect who will exist.

\paragraph{Just savings.} Rawls held that each generation must save for its successors at a rate that members of each generation, not knowing to which they belong, would agree to \citep[\S 44]{Rawls1971}. The principle is formulated for generations, not persons, and it is in our sense impersonal: it concerns what is owed to the occupants of later generational positions, whoever they are. Our argument explains why intergenerational justice has had to be formulated in this way, and why attempts to formulate it as a set of claims held by particular future persons have met with difficulties.

\section{Objections and replies}\label{s:objections}

\begin{objection}[Every wrong has a victim]
Wrongness is relational: to act wrongly is to wrong someone. An impersonal obligation, whose violation wrongs no one, is not a moral obligation at all, but at most a reason to promote value.
\end{objection}

\begin{reply}
The objection begs the question. There are uncontroversial examples of wrongs without particular victims: the destruction of a great work of art that no one will ever see, or the waste of goods that would have benefited someone or other, though no one in particular. The relational conception of wrongness is plausible for the core of interpersonal morality, and we do not deny it there; we deny that it exhausts morality. The non-identity problem is a demonstration that it does not: Depletion is wrong, and no one is wronged. An account of moral obligation that cannot accommodate this has to deny Principle~\ref{p:depletion}, which is the bullet that Narveson and Boonin bite.
\end{reply}

\begin{objection}[The claim condition is ad hoc]
The second part of Principle~\ref{p:claim}, which requires that some alternative honor the claim with its holder existing, was introduced to exclude the claims of future people. It has no independent support.
\end{objection}

\begin{reply}
It is supported by what claims are for. A claim gives its holder standing to demand a certain treatment and to complain if it is withheld. Standing to demand presupposes that the demanded treatment is something the agent could give the claimant. Consider a person who claims that he ought to have been born to richer parents, whose having been born to richer parents would have made him a different person. We do not regard his poorer parents as having infringed a claim of his, however much we regret his circumstances; we regard the complaint as confused. The second part of Principle~\ref{p:claim} states the condition whose violation makes such complaints confused. It applies equally to cases with no intergenerational element, and it does not exclude the claims of future people who will exist whatever we do, such as the child in the buried-glass case.
\end{reply}

\begin{objection}[Future people will resent us]
The people of the twenty-fourth century will resent the generation that depleted their resources, and their resentment will be appropriate. Appropriate resentment is a mark of having been wronged.
\end{objection}

\begin{reply}
Their indignation will be appropriate: the depleting generation acted wrongly, and indignation is the appropriate response to wrongdoing, whether or not one is its victim. Resentment, the response of a victim to being wronged, is not appropriate, since none of them would exist had the generation acted otherwise. The distinction between resentment and indignation is the distinction between directed and impersonal obligation, and the case confirms our view that the moral responses available to future people are those of third parties to a wrong. These responses are often strong; impersonal obligations are not weak obligations.
\end{reply}

\begin{objection}[Collective holders]
The obligation is owed not to any future person but to future generations as a collective, which exists under either option and is made worse off by depletion.
\end{objection}

\begin{reply}
Whether a future generation exists under both options depends on how generations are individuated. If a generation is individuated by its members, it does not exist under both options, since its members differ. If it is individuated by a temporal position, the generation of the twenty-third century, then it is not a holder of claims but a position, and an obligation owed to a position is owed to whoever occupies it, which is to say to no one in particular. The collective proposal is therefore either false or a restatement of our thesis.
\end{reply}

\section{What the argument does not establish}\label{s:limits}

The argument establishes that the obligations to future generations as such are impersonal, given the comparative condition on harming, the claim condition on wronging, and the datum that depletion is wrong. It does not establish the content of these obligations in detail: how much a generation must save, or what level of emissions it may allow, depends on empirical questions and on the deontic principles that link impersonal value to obligation. It does not establish that all obligations concerning future people are impersonal, and the buried-glass case shows that some are not. It does not settle the question whether coming into existence can benefit a person, a question that our argument does not need. And it does not establish the population axiology on which the obligations are grounded; for that we rely on ETH-09.

\section{Conclusion}

We owe something to future generations, but not to anyone in particular. Policies that affect the conditions of life of future generations also affect who will exist, and those who exist as a result of a bad policy are neither made worse off by it nor holders of claims that the policy could have honored. Since such policies are nonetheless wrong, the obligations they violate are impersonal. The attempts to find a particular claimant either fail in the central cases or relocate the claim in a type, a position or a description, which concedes the point. The obligations are grounded in the impersonal value of outcomes, and their force is undiminished by their having no one to whom they are owed.

\begin{thebibliography}{99}\small
\bibitem[Adams(1979)]{Adams1979} Adams, R.~M. (1979). Existence, self-interest, and the problem of evil. \emph{No{\^u}s}, 13(1), 53--65.
\bibitem[Boonin(2014)]{Boonin2014} Boonin, D. (2014). \emph{The Non-Identity Problem and the Ethics of Future People}. Oxford University Press.
\bibitem[Broome(2012)]{Broome2012} Broome, J. (2012). \emph{Climate Matters: Ethics in a Warming World}. W.~W. Norton.
\bibitem[Gardiner(2011)]{Gardiner2011} Gardiner, S.~M. (2011). \emph{A Perfect Moral Storm: The Ethical Tragedy of Climate Change}. Oxford University Press.
\bibitem[Hanser(1990)]{Hanser1990} Hanser, M. (1990). Harming future people. \emph{Philosophy \& Public Affairs}, 19(1), 47--70.
\bibitem[Hare(2007)]{Hare2007} Hare, C. (2007). Voices from another world: Must we respect the interests of people who do not, and will never, exist? \emph{Ethics}, 117(3), 498--523.
\bibitem[Harman(2004)]{Harman2004} Harman, E. (2004). Can we harm and benefit in creating? \emph{Philosophical Perspectives}, 18, 89--113.
\bibitem[Heyd(1992)]{Heyd1992} Heyd, D. (1992). \emph{Genethics: Moral Issues in the Creation of People}. University of California Press.
\bibitem[Kavka(1982)]{Kavka1982} Kavka, G.~S. (1982). The paradox of future individuals. \emph{Philosophy \& Public Affairs}, 11(2), 93--112.
\bibitem[Kripke(1980)]{Kripke1980} Kripke, S.~A. (1980). \emph{Naming and Necessity}. Harvard University Press.
\bibitem[Kumar(2003)]{Kumar2003} Kumar, R. (2003). Who can be wronged? \emph{Philosophy \& Public Affairs}, 31(2), 99--118.
\bibitem[Narveson(1967)]{Narveson1967} Narveson, J. (1967). Utilitarianism and new generations. \emph{Mind}, 76(301), 62--72.
\bibitem[Parfit(1984)]{Parfit1984} Parfit, D. (1984). \emph{Reasons and Persons}. Clarendon Press.
\bibitem[Rawls(1971)]{Rawls1971} Rawls, J. (1971). \emph{A Theory of Justice}. Harvard University Press.
\bibitem[Roberts(2010)]{Roberts2010} Roberts, M.~A. (2010). \emph{Abortion and the Moral Significance of Merely Possible Persons: Finding Middle Ground in Hard Cases}. Springer.
\bibitem[Schwartz(1978)]{Schwartz1978} Schwartz, T. (1978). Obligations to posterity. In R.~I. Sikora and B. Barry (eds.), \emph{Obligations to Future Generations}. Temple University Press.
\bibitem[Shiffrin(1999)]{Shiffrin1999} Shiffrin, S.~V. (1999). Wrongful life, procreative responsibility, and the significance of harm. \emph{Legal Theory}, 5(2), 117--148.
\bibitem[Woodward(1986)]{Woodward1986} Woodward, J. (1986). The non-identity problem. \emph{Ethics}, 96(4), 804--831.
\end{thebibliography}

\end{document}
